\subsection{The countable form of the primary separation test}

The following is a rechecked consequence of the supplied REVLOC G.7,
not a newly discovered reduction of the open problem.  We include a
second compactness proof from an opened original source, so that the
countable bound does not depend on an unexamined inverse-limit or
phantom-map assertion.

\begin{claim}{A4.2}{Proved}{Countable primary sphere-null spaces suffice}
\label{A42}
Let
\[
 \mathcal A_* = \Nul_2\cap
 \{X:\widetilde H_*(X;\mathbb Z[1/2])=0\}.
\]
Every countable pointed subcomplex of a CW complex $X\in\mathcal A_*$
is contained in a countable pointed subcomplex $Y\in\mathcal A_*$.
Consequently $W_*=W_S\cap W_{\mathbb Z[1/2]}$ is accessible in ZFC.
Moreover
\[
 W_*=\min\operatorname{Rev}(BC_2\to *)
\]
if and only if every \emph{countable} connected pointed CW complex $X$
with $2$-primary reduced integral homology and
$\Map_*(X,S^n_{(2)})\simeq *$ for all $n\geq2$ satisfies
\[
                       P_{BC_2}(\Sigma X)\simeq *.
 \tag{A4.2.1}
\]
Thus, if the first inclusion of A4.1 is strict, it has a countable
aspherical-space witness.  This statement decides neither that
strictness nor the size of the full fixed-stable-part fibre.
\end{claim}
\begin{proof}
Write $Z_n=(S^n)^\wedge_2$.  The precise classical input is
\citet[Section~3, pp.~34--41, Property~iii on p.~49, and the
prime-wise version on p.~52]{A4SullivanGenetics}.  For each $n$ there
is a countable inverse system of spaces $T_{n,j}$ with finite
$2$-primary homotopy groups, such that, naturally for \emph{every}
pointed CW complex $V$,
\[
 [V,Z_n]_*
   \cong\lim_j[V,T_{n,j}]_*
   \cong\lim_{j,K\subset V\text{ finite}}[K,T_{n,j}]_*.
 \tag{A4.2.2}
\]
The last sets are finite.  The first equality is Sullivan's
\emph{represented inverse limit of compact Brownian functors}; it
is not an interchange of $\pi_0$ with an arbitrary homotopy limit.
His finite-source countability hypothesis holds for $S^n$; equivalently,
one can use the countably many truncated finite $2$-homotopy types
and their finitely many homotopy classes of maps from $S^n$.
The finite-subcomplex formula is Property~I and Example~II on
pp.~36--37.  Both formulas also apply to $\Sigma^sV$, for all $s\geq0$.

We spell out the compactness consequence that is needed.  If
$[\Sigma^sX,Z_n]_*=*$, then for every finite pointed $K\subset X$
and every index $j$ there are a finite pointed $L\supset K$ in $X$
and a later index $t$ such that the image of
\[
 [\Sigma^sL,T_{n,t}]_*\longrightarrow
 [\Sigma^sK,T_{n,j}]_*
 \tag{A4.2.3}
\]
is the singleton consisting of the constant map.  Indeed the
inverse system of finite pointed sets in \textup{(A4.2.2)} has
singleton limit.  Its image in any fixed coordinate is the
intersection of the images of later coordinates: a compatible
family with prescribed coordinate exists whenever every finite
collection of the compatibility conditions can be solved, by
compactness of the product of finite discrete sets.  Since the fixed
coordinate is finite, finitely many later images already have the
same singleton intersection.  A common later coordinate gives
\textup{(A4.2.3)}.  No assertion that $[X,T_{n,j}]_*=*$ is made.

For a rationally acyclic source $V$, the arithmetic square of the
$2$-local sphere gives
\[
 \Map_*(V,S^n_{(2)})\simeq\Map_*(V,Z_n).
 \tag{A4.2.4}
\]
Here both rational vertices are simply connected rational spaces,
hence rational-homology-local, so their pointed mapping spaces from
$V$ are contractible.  For completeness, the arithmetic square can
be checked on homotopy groups: these groups of the integral sphere
are finitely generated, and localization and the completion
calculation in Sullivan's Theorem~3.1 give the exact sequences
\[
 0\longrightarrow A\otimes\mathbb Z_{(2)}
 \longrightarrow (A\otimes\mathbb Z_2)\oplus(A\otimes\mathbb Q)
 \longrightarrow A\otimes\mathbb Q_2\longrightarrow0
\]
for each such group $A$.  The homotopy pullback therefore has exactly
the homotopy groups of $S^n_{(2)}$.  This checks the target square
without a finiteness assumption on $V$.

Start with the specified countable subcomplex of $X$.  At each stage
adjoin the finite subcomplexes $L$ from \textup{(A4.2.3)} for every
finite $K$ in the current subcomplex and every $n,j,s$.  These are
countably many choices.  At the same stage, for every reduced cellular
cycle with coefficients in the countable ring $\mathbb Z[1/2]$ in
the current subcomplex, adjoin the finite support of one bounding
chain in $X$, together with its cell closures.  Such chains exist
because $X$ is $\mathbb Z[1/2]$-acyclic.  Iterate countably many times
and call the union $Y$.  Every reduced cycle in $Y$ occurs at a finite
stage and bounds at the next, so $Y$ is $\mathbb Z[1/2]$-acyclic;
reduced degree zero also proves connectedness.

Every map $\Sigma^sY\to Z_n$ has zero image in each coordinate
$[\Sigma^sK,T_{n,j}]_*$, by the corresponding witness
\textup{(A4.2.3)} lying in $Y$.  Formula \textup{(A4.2.2)} says the
map itself is null.  This includes $s=0$, so the mapping space has
one component, and all $s>0$ make its homotopy groups zero.  Thus
$Y$ is completed-sphere-null.  Its rational acyclicity and
\textup{(A4.2.4)} put it in $\mathcal A_*$.  This proves the
enlargement assertion.

The bounded-subcomplex argument of A.7 now applies with
$\mathcal A=\mathcal A_*$ and proves accessibility.  In particular
all members of $\mathcal A_*$ are filtered homotopy colimits of
their countable $\mathcal A_*$-subcomplexes: the enlargement property
makes that poset filtered and its union is the original complex.
The lower reversible aspherical class is closed under these
homotopy colimits.  It therefore contains $\mathcal A_*$ exactly
when it contains every countable member.  Finally A34.2 identifies
membership in the lower class with \textup{(A4.2.1)}, including its
one required initial suspension.
\end{proof}
